\section{Keccak reducido y modelo de actividad}
\label{sec:keccak_reducido}

Keccak es una familia de funciones criptográficas basada en una permutación
sobre un estado tridimensional de bits. El estándar SHA-3 utiliza
permutaciones de la familia Keccak-$p$ y define cada ronda mediante la
composición de cinco transformaciones:

\[
\theta
\longrightarrow
\rho
\longrightarrow
\pi
\longrightarrow
\chi
\longrightarrow
\iota.
\]

Las transformaciones $\theta$, $\rho$ y $\pi$ conforman la parte lineal de
la ronda. La transformación $\chi$ introduce la no linealidad, mientras que
$\iota$ incorpora una constante dependiente del número de ronda
\cite{nist_fips202,bertoni2011keccak_reference}.

En este trabajo se conserva la estructura de la ronda de Keccak, pero se
emplean tamaños de palabra reducidos para hacer viable la construcción y
resolución del modelo MILP.


\subsection{Representación del estado}
\label{subsec:representacion_estado}

El estado al inicio de la ronda $r$ se representa mediante el arreglo
binario

\[
A_r[x,y,k] \in \{0,1\},
\]

donde

\[
x,y \in \{0,1,2,3,4\},
\qquad
k \in \{0,\ldots,z-1\}.
\]

Las coordenadas $(x,y)$ identifican una de las 25 palabras o
\emph{lanes} del estado, mientras que $k$ identifica una posición dentro
de una palabra de longitud $z$. Por tanto, el tamaño total del estado es

\[
25z \text{ bits}.
\]

Se analizan los tamaños de palabra $z=4$ y $z=8$. La
Tabla~\ref{tab:configuraciones_keccak} resume las configuraciones
consideradas.

\begin{table}[htbp]
    \centering
    \caption{Configuraciones reducidas de Keccak consideradas en el trabajo.}
    \label{tab:configuraciones_keccak}
    \begin{tabular}{cccc}
        \hline
        $z$ & Tamaño del estado & Palabras & S-boxes por ronda \\
        \hline
        4 & 100 bits & 25 & 20 \\
        8 & 200 bits & 25 & 40 \\
        \hline
    \end{tabular}
\end{table}

Cada aplicación local de $\chi$ procesa cinco bits que comparten las
coordenadas $(y,k)$ y recorren los cinco valores posibles de $x$. Como
existen cinco valores de $y$ y $z$ posiciones dentro de cada palabra, el
número de aplicaciones locales de $\chi$ por ronda es

\[
5z.
\]


\subsection{Transformaciones de una ronda}
\label{subsec:transformaciones_ronda}

Para distinguir los estados intermedios de una ronda se utiliza la siguiente
notación:

\[
T_r = \theta(A_r),
\]

\[
R_r = \rho(T_r),
\]

\[
B_r = \pi(R_r),
\]

\[
C_r = \chi(B_r),
\]

\[
A_{r+1} = \iota(C_r,r).
\]

El estado $B_r$ corresponde a la entrada de la capa no lineal $\chi$ y
será utilizado posteriormente para determinar qué S-boxes se encuentran
activas.


\subsubsection{Transformación \texorpdfstring{$\theta$}{theta}}

La transformación $\theta$ introduce difusión entre las columnas del
estado. Primero se calcula la paridad de cada columna:

\[
P_r[x,k]
=
\bigoplus_{y=0}^{4}
A_r[x,y,k].
\]

A partir de estas paridades se define el término de difusión

\[
D_r[x,k]
=
P_r[(x-1)\bmod 5,k]
\oplus
P_r[(x+1)\bmod 5,(k-1)\bmod z].
\]

La salida de $\theta$ es

\[
T_r[x,y,k]
=
A_r[x,y,k]
\oplus
D_r[x,k].
\]

La transformación $\theta$ es lineal sobre $\mathrm{GF}(2)$ y distribuye
la información de la paridad de una columna hacia columnas vecinas.


\subsubsection{Transformaciones \texorpdfstring{$\rho$}{rho} y
\texorpdfstring{$\pi$}{pi}}

La transformación $\rho$ rota cada palabra según un desplazamiento
predefinido. Los desplazamientos especificados por Keccak se interpretan
módulo $z$:

\[
R_r[x,y,k]
=
T_r
\left[
x,
y,
\bigl(k-r[x,y]\bigr)\bmod z
\right],
\]

donde $r[x,y]$ representa el desplazamiento asignado a la palabra
$(x,y)$.

La transformación $\pi$ reorganiza las palabras dentro del estado:

\[
B_r
\left[
y,
(2x+3y)\bmod 5,
k
\right]
=
R_r[x,y,k].
\]

Las transformaciones $\rho$ y $\pi$ son permutaciones de posiciones. Por
tanto, no modifican los valores de los bits, sino únicamente su ubicación
dentro del estado.


\subsubsection{Transformación \texorpdfstring{$\chi$}{chi}}

La transformación $\chi$ constituye la única capa no lineal de la ronda.
Para cada posición $(y,k)$ opera sobre los cinco bits obtenidos al variar
$x$:

\[
C_r[x,y,k]
=
B_r[x,y,k]
\oplus
\left[
\left(
1 \oplus B_r[(x+1)\bmod 5,y,k]
\right)
B_r[(x+2)\bmod 5,y,k]
\right].
\]

Cada conjunto

\[
\left(
B_r[0,y,k],
B_r[1,y,k],
B_r[2,y,k],
B_r[3,y,k],
B_r[4,y,k]
\right)
\]

se interpreta en este trabajo como una S-box local de cinco bits.

Esta denominación no implica que Keccak utilice una tabla de sustitución
independiente en el sentido tradicional. Se emplea porque cada grupo local
de cinco bits constituye una unidad no lineal cuya actividad puede
contabilizarse durante la propagación diferencial.


\subsubsection{Transformación \texorpdfstring{$\iota$}{iota}}

La transformación $\iota$ incorpora una constante de ronda en la palabra
ubicada en las coordenadas $(0,0)$. Su acción puede expresarse como

\[
A_{r+1}[x,y,k]
=
\begin{cases}
C_r[0,0,k] \oplus RC_r[k],
    & \text{si } (x,y)=(0,0),\\[4pt]
C_r[x,y,k],
    & \text{en otro caso},
\end{cases}
\]

donde $RC_r[k]$ representa el bit correspondiente de la constante de la
ronda $r$.

Para los tamaños de palabra $z=4$ y $z=8$, las constantes se restringen a
las primeras $z$ posiciones de la palabra.


\subsection{Modelo diferencial pareado}
\label{subsec:modelo_diferencial_pareado}

El análisis considera dos ejecuciones simultáneas de Keccak, denominadas
ejecución izquierda y ejecución derecha. Sus estados al inicio de la ronda
$r$ se representan mediante

\[
A_r^L
\qquad \text{y} \qquad
A_r^R.
\]

La diferencia entre ambos estados se define mediante XOR:

\[
\Delta A_r
=
A_r^L \oplus A_r^R.
\]

De forma análoga, los estados inmediatamente anteriores a la capa $\chi$
son

\[
B_r^L
=
(\pi \circ \rho \circ \theta)(A_r^L),
\]

\[
B_r^R
=
(\pi \circ \rho \circ \theta)(A_r^R),
\]

y su diferencia es

\[
\Delta B_r
=
B_r^L \oplus B_r^R.
\]

La utilización de dos ejecuciones concretas permite representar
exactamente el comportamiento de la capa no lineal. En lugar de aproximar
la propagación diferencial de $\chi$, se evalúa la transformación sobre
ambos estados y posteriormente se calcula la diferencia entre sus salidas.

La constante de ronda aplicada por $\iota$ es idéntica en las dos
ejecuciones. En consecuencia, se cancela al calcular la diferencia:

\[
\left(
C_r^L \oplus RC_r
\right)
\oplus
\left(
C_r^R \oplus RC_r
\right)
=
C_r^L \oplus C_r^R.
\]

Aunque $\iota$ no modifica la diferencia, se conserva en la implementación
para mantener una representación exacta de los estados absolutos y permitir
la reconstrucción completa de los testigos.


\subsection{Definición de S-box activa}
\label{subsec:sbox_activa}

Una aplicación local de $\chi$ en la posición $(y,k)$ de la ronda $r$ se
considera activa cuando al menos uno de sus cinco bits de entrada presenta
una diferencia no nula.

Se introduce el indicador binario

\[
a_{r,y,k} \in \{0,1\},
\]

definido por

\[
a_{r,y,k}
=
\bigvee_{x=0}^{4}
\Delta B_r[x,y,k].
\]

Por tanto,

\[
a_{r,y,k}=1
\]

si existe al menos un valor de $x$ para el cual

\[
\Delta B_r[x,y,k]=1.
\]

El número de S-boxes activas en la ronda $r$ se define como

\[
m_r
=
\sum_{y=0}^{4}
\sum_{k=0}^{z-1}
a_{r,y,k}.
\]

Para una ejecución de $R$ rondas, la actividad total es

\[
N_{\mathrm{act}}
=
\sum_{r=0}^{R-1}
m_r
=
\sum_{r=0}^{R-1}
\sum_{y=0}^{4}
\sum_{k=0}^{z-1}
a_{r,y,k}.
\]

El problema de optimización estudiado consiste en minimizar
$N_{\mathrm{act}}$ bajo la condición de que la diferencia inicial sea no
nula.


\subsection{Propiedades estructurales utilizadas}
\label{subsec:propiedades_estructurales}

La interpretación de los resultados utiliza las siguientes propiedades de
Keccak:

\begin{enumerate}
    \item Las transformaciones $\theta$, $\rho$ y $\pi$ son biyectivas.
    Por tanto, una diferencia no nula al inicio de una ronda no puede
    convertirse en una diferencia nula antes de alcanzar la capa $\chi$:

    \[
    \Delta A_r \neq 0
    \Longrightarrow
    \Delta B_r \neq 0.
    \]

    \item La ronda completa de Keccak es una permutación. En consecuencia,
    dos estados distintos continúan siendo distintos después de aplicar
    cualquier número de rondas:

    \[
    A_r^L \neq A_r^R
    \Longrightarrow
    A_{r+1}^L \neq A_{r+1}^R.
    \]

    \item La transformación $\iota$ no modifica la diferencia entre las
    dos ejecuciones porque aplica la misma constante a ambas.

    \item Una diferencia no nula en $B_r$ activa necesariamente al menos
    una aplicación local de $\chi$:

    \[
    \Delta B_r \neq 0
    \Longrightarrow
    m_r \geq 1.
    \]
\end{enumerate}

Estas propiedades permiten establecer cotas inferiores básicas para una o
varias rondas. La forma en que se incorporan al modelo MILP y se utilizan
para certificar los resultados se desarrolla en las secciones siguientes.